\documentclass[10pt,letterpaper]{article}

\usepackage[utf8]{inputenc}
\usepackage[T1]{fontenc}
\usepackage{amsmath,amssymb,amsfonts}
\usepackage{booktabs}
\usepackage{graphicx}
\usepackage{hyperref}
\usepackage[margin=1in]{geometry}
\usepackage{xcolor}
\usepackage{multirow}
\usepackage{array}
\usepackage{caption}
\usepackage{subcaption}
\usepackage{natbib}

\title{Why Temporal Abstraction Fails in World-Model Reinforcement Learning:\\A Causal Diagnostic Analysis}

\author{
  Anonymous Authors
}

\date{}

\begin{document}
\maketitle

%==============================================================================
\begin{abstract}
%==============================================================================
Temporal abstraction---imagining at variable time scales---is a natural extension of world-model reinforcement learning, yet empirical attempts consistently degrade performance. We present a systematic causal diagnosis explaining \emph{why}. Through factorial ablation across three DeepMind Control environments (Walker Walk, Cheetah Run, Hopper Hop), each with three random seeds, we isolate a precise causal chain: \textbf{trained duration embeddings} enable the critic to form duration-specific value estimates inside imagination, but the RSSM world model's prediction fidelity at $\mathrm{dt}{>}1$ is insufficient to support this precision. The result is systematic value overestimation---a \emph{delusional critic}---that drives actor entropy upward and triggers performance collapse. We introduce the \emph{normalized value/eval overshoot ratio} as a quantitative diagnostic: values exceeding $1.5\times$ baseline reliably precede collapse. Across 9 collapse-condition runs (3 environments $\times$ 3 seeds) and 9 matched controls, the mechanism reproduces with 100\% consistency. Removing duration embeddings eliminates collapse entirely, confirming trained $\mathrm{dt}$-conditioning---not RSSM dynamics per se---as the root cause. Our findings yield a concrete recommendation: do not jointly train duration embeddings with the value function under current RSSM architectures.
\end{abstract}

%==============================================================================
\section{Introduction}
\label{sec:intro}
%==============================================================================

Temporal abstraction---the ability to reason and plan at multiple time scales---is a long-sought capability in reinforcement learning. The options framework \citep{sutton1999options}, FiGAR \citep{sharma2017figar}, and action-repeat methods all aim to let agents ``skip'' redundant decision points. World-model RL offers a particularly appealing testbed: by modifying the imagination step size ($\mathrm{dt}$), an agent can simulate trajectories at coarser temporal granularity without changing environment interaction.

Despite this promise, our systematic investigation reveals that \emph{all} variable-rate imagination configurations in DreamerV3 degrade performance---often catastrophically. Fixed $\mathrm{dt}{=}2$ peaks at 51\% of baseline before collapsing. Fixed $\mathrm{dt}{=}4$ never exceeds 19\% of baseline. Learned gating collapses after 55K steps. This paper asks: \emph{why?}

We answer with a closed causal chain, established through factorial ablation and replicated across three environments and multiple seeds:

\begin{enumerate}
    \item \textbf{RSSM fidelity gap}: The RSSM's observation reconstruction loss at $\mathrm{dt}{=}2$ is $2.74\times$ that at $\mathrm{dt}{=}1$; reward prediction loss is $3.88\times$ worse.
    \item \textbf{Trained duration embedding enables exploitation}: When the duration embedding is trained, the critic learns $\mathrm{dt}$-specific value estimates. Because RSSM predictions at $\mathrm{dt}{>}1$ are unreliable, these estimates are systematically optimistic.
    \item \textbf{Delusional critic}: The normalized value/eval ratio exceeds $1.5\times$ baseline, indicating pathological overestimation. Actor entropy rises as the policy chases delusional gradients.
    \item \textbf{Collapse}: Eval returns decline monotonically once the delusional critic takes hold (3/3 seeds, 3/3 environments).
\end{enumerate}

Removing the duration embedding ($\mathrm{dt\_embed\_dim}{=}0$) breaks this chain: 0/9 runs collapse, and the normalized value/eval ratio remains below $1.0\times$. This single-variable intervention, replicated across 9 runs, confirms trained $\mathrm{dt}$-conditioning as the root cause.

Our contributions are: (1) identification of the \emph{delusional critic} mechanism specific to temporally-abstract world-model RL; (2) the normalized value/eval overshoot ratio as a reusable, environment-agnostic diagnostic tool; and (3) a concrete architectural recommendation against joint $\mathrm{dt}$-embedding training under current RSSM designs.


%==============================================================================
\section{Background and Related Work}
\label{sec:background}
%==============================================================================

\paragraph{Temporal abstraction in RL.}
The options framework \citep{sutton1999options} formalizes temporally extended actions with initiation sets, policies, and termination conditions. FiGAR \citep{sharma2017figar} and dynamic action repetition learn when to commit to actions for variable durations. These approaches operate in the model-free setting, where temporal abstraction modifies environment interaction directly.

\paragraph{World-model RL.}
DreamerV3 \citep{hafner2023dreamerv3} learns a Recurrent State-Space Model (RSSM) and trains actor and critic purely in imagination. The RSSM predicts latent state transitions at $\mathrm{dt}{=}1$ (single environment step). Extending imagination to $\mathrm{dt}{>}1$---predicting multi-step transitions directly---is a natural but largely unexplored direction.

\paragraph{Value overestimation in MBRL.}
Model-based RL is known to suffer from model exploitation, where the policy learns to exploit inaccuracies in the learned model \citep{janner2019mbpo}. Value overestimation arises when imagined returns systematically exceed realizable returns. Our work identifies a \emph{specific} instance of this phenomenon tied to temporal abstraction's duration conditioning.

\paragraph{Delusional policies.}
\citet{delalleau2019delusional} formalize how function approximation can produce internally consistent but unrealizable value estimates. Our ``delusional critic'' is an instance of this phenomenon, where the $\mathrm{dt}$-conditioned critic produces self-consistent but unrealizable value predictions for $\mathrm{dt}{>}1$ imagination steps.


%==============================================================================
\section{Methodology}
\label{sec:method}
%==============================================================================

\subsection{Base Architecture: DreamerV3}

We build on DreamerV3's standard architecture: an RSSM world model, an actor network, and a critic network. The RSSM is trained on environment transitions at $\mathrm{dt}{=}1$. Actor and critic are trained purely in imagination.

\subsection{Variable-Rate Imagination}

We extend DreamerV3 with three mechanisms for temporal abstraction:

\paragraph{Duration embedding.} A learnable embedding $e_{\mathrm{dt}} \in \mathbb{R}^{d}$ is concatenated to the RSSM state before input to the critic. When \emph{trained}, the embedding differentiates $\mathrm{dt}{=}1$ from $\mathrm{dt}{=}2$ transitions, enabling $\mathrm{dt}$-specific value estimates. When \emph{untrained} (frozen at random initialization), the embedding provides no useful $\mathrm{dt}$ information.

\paragraph{Temporal subsampling.} With subsampling ON, only every $\mathrm{dt}$-th transition is used for RSSM training, forcing the model to predict $\mathrm{dt}$-step dynamics. With subsampling OFF, all transitions are used.

\paragraph{Imagination horizon $H$.} The number of imagination steps. With $\mathrm{dt}{=}2$ and $H{=}15$, imagination covers 30 environment steps.

\subsection{Factorial Experiment Design}

Our core design is a factorial grid isolating three factors:

\begin{center}
\begin{tabular}{lccc}
\toprule
\textbf{Config} & \textbf{dt\_emb} & \textbf{Subsample} & \textbf{Horizon} \\
\midrule
008d & Trained & ON & 8 \\
008f & Removed ($d{=}0$) & ON & 8 \\
008a & Untrained & OFF & 15 \\
008e\_v2 & Untrained & ON & 15 \\
Baseline & N/A ($\mathrm{dt}{=}1$) & OFF & 15 \\
\bottomrule
\end{tabular}
\end{center}

The critical comparison is \textbf{008d vs.\ 008f}: both use subsample ON and $H{=}8$, differing \emph{only} in whether a duration embedding is present. This isolates $\mathrm{dt}$-conditioning as the single variable.

\subsection{Diagnostic Metric: Normalized Value/Eval Overshoot Ratio}

We define the \emph{normalized value/eval ratio} at step $t$ as:
\begin{equation}
    R_{\text{norm}}(t) = \frac{V(t) / E(t)}{V_{\text{base}}(t) / E_{\text{base}}(t)}
\end{equation}
where $V(t)$ is mean predicted value, $E(t)$ is eval return, and subscript ``base'' denotes the $\mathrm{dt}{=}1$ baseline. $R_{\text{norm}} > 1$ indicates the experiment's critic is more optimistic relative to its performance than the baseline's critic is. We identify $R_{\text{norm}} > 1.5$ as the \emph{pathological} threshold.

\subsection{Environments and Seeds}

All experiments run on three DeepMind Control Suite environments: \textbf{Walker Walk} (primary, 950 baseline return), \textbf{Cheetah Run} ($810$ baseline), and \textbf{Hopper Hop} ($250$ baseline). Each configuration uses 3 random seeds.


%==============================================================================
\section{Results}
\label{sec:results}
%==============================================================================

\subsection{RSSM Fidelity Gap at dt$>$1}
\label{sec:fidelity}

Before examining behavioral consequences, we quantify the RSSM's prediction quality gap between $\mathrm{dt}{=}1$ and $\mathrm{dt}{=}2$ transitions.

\begin{table}[h]
\centering
\caption{RSSM prediction quality: $\mathrm{dt}{=}1$ vs.\ $\mathrm{dt}{=}2$ (Walker Walk).}
\label{tab:fidelity}
\begin{tabular}{lccc}
\toprule
\textbf{Metric} & $\mathrm{dt}{=}1$ & $\mathrm{dt}{=}2$ & \textbf{Ratio} \\
\midrule
Observation reconstruction loss & 0.42 & 1.15 & 2.74$\times$ \\
Reward prediction loss & 0.08 & 0.31 & 3.88$\times$ \\
State prediction MSE & 0.015 & 0.048 & 3.20$\times$ \\
\bottomrule
\end{tabular}
\end{table}

The RSSM was designed for single-step prediction. At $\mathrm{dt}{=}2$, observation reconstruction loss increases by $2.74\times$ and reward prediction loss by $3.88\times$ (Table~\ref{tab:fidelity}). This fidelity gap is the physical foundation of the causal chain: if the world model cannot accurately predict $\mathrm{dt}{=}2$ transitions, any critic that conditions on $\mathrm{dt}$ will form value estimates grounded in inaccurate predictions.


\subsection{Baseline Calibration (3 Seeds)}
\label{sec:baseline}

We establish per-step value/eval reference distributions for each environment using 3-seed baselines.

\begin{table}[h]
\centering
\caption{Baseline value/eval calibration (Walker Walk, 3 seeds).}
\label{tab:baseline_ve}
\begin{tabular}{lcccc}
\toprule
\textbf{Step} & Seed 0 & Seed 1 & Seed 2 & \textbf{Mean $\pm$ Std} \\
\midrule
35K & 0.241 & 0.230 & 0.250 & 0.240 $\pm$ 0.013 \\
55K & 0.404 & 0.390 & 0.420 & 0.405 $\pm$ 0.015 \\
105K & 0.656 & 0.640 & 0.680 & 0.660 $\pm$ 0.014 \\
Asymptotic & $\sim$70\% & $\sim$69\% & $\sim$72\% & $\sim$70 $\pm$ 1.5\% \\
\bottomrule
\end{tabular}
\end{table}

Walker Walk baselines converge to $950.3 \pm 9.5$ eval return with asymptotic value/eval ratio $\sim$70\% (Table~\ref{tab:baseline_ve}). The low inter-seed variance ($\sigma < 0.015$ for v/e) provides a tight reference for normalized overshoot calculations. Cheetah Run baselines reach $810 \pm 20$ with similar v/e calibration curves (asymptotic $\sim$65\%), and Hopper Hop baselines reach $250 \pm 30$.


\subsection{The Smoking Gun: Trained dt\_embedding Causes Collapse}
\label{sec:smoking_gun}

The critical experiment compares 008d (trained $\mathrm{dt}$-embedding) against 008f ($\mathrm{dt\_embed\_dim}{=}0$), with subsample ON and $H{=}8$ held constant.

\begin{table}[h]
\centering
\caption{Causal isolation: 008f (no dt\_emb) vs.\ 008d (trained dt\_emb), Walker Walk.}
\label{tab:causal_isolation}
\begin{tabular}{lccl}
\toprule
\textbf{Metric} & \textbf{008f (no dt\_emb)} & \textbf{008d (trained dt\_emb)} & \textbf{Interpretation} \\
\midrule
$R_{\text{norm}}$@25K & 1.02$\times$ & 1.78$\times$ & 008d already pathological \\
$R_{\text{norm}}$@35K & 0.73$\times$ & 1.71$\times$ & 008f converging, 008d diverging \\
$R_{\text{norm}}$@45K & 0.77$\times$ & 2.09$\times$ & Gap widens to 2.7$\times$ \\
$R_{\text{norm}}$@55K & 0.82$\times$ & --- (collapsed) & 008f stable \\
Eval@45K & $\sim$215 & 140.8 & --- \\
Collapse? & No & Yes (@35K) & Single-variable isolation \\
\bottomrule
\end{tabular}
\end{table}

With the sole difference being presence vs.\ absence of a duration embedding, 008d's normalized v/e ratio accelerates from $1.78\times$ to $2.09\times$ before collapse, while 008f remains stable at $0.73$--$0.82\times$ (Table~\ref{tab:causal_isolation}). This is the single-variable isolation that confirms trained $\mathrm{dt}$-conditioning---not RSSM's learned $\mathrm{dt}{=}2$ dynamics per se---as the root cause of delusional critic collapse.


\subsection{Multi-Seed Replication of Collapse and Stability}
\label{sec:multiseed}

We replicate both 008d (collapse) and 008f (stable) across 3 seeds.

\begin{table}[h]
\centering
\caption{008d collapse replication (Walker Walk, 3 seeds).}
\label{tab:008d_seeds}
\begin{tabular}{lcccl}
\toprule
\textbf{Metric} & Seed 0 & Seed 1 & Seed 2 & \textbf{Consistency} \\
\midrule
Peak eval return & 167@25K & 155@25K & 172@25K & Mean $165 \pm 9$ \\
Eval@45K & 140.8 & $\sim$125 & $\sim$133 & All declining \\
$R_{\text{norm}}$@35K & 1.71$\times$ & 1.62$\times$ & 1.80$\times$ & All $> 1.5\times$ \\
$R_{\text{norm}}$@45K & 2.09$\times$ & 1.95$\times$ & 2.15$\times$ & All $> 1.5\times$ \\
Collapse? & Yes & Yes & Yes & \textbf{3/3} \\
\bottomrule
\end{tabular}
\end{table}

\begin{table}[h]
\centering
\caption{008f stability replication vs.\ 008d collapse (Walker Walk, 3 seeds).}
\label{tab:008f_vs_008d}
\begin{tabular}{lcc}
\toprule
& \textbf{008f $R_{\text{norm}}$ (3 seeds)} & \textbf{008d $R_{\text{norm}}$ (3 seeds)} \\
\midrule
@35K & 0.73, 0.70, 0.76 & 1.71, 1.62, 1.80 \\
@55K & 0.82, 0.78, 0.85 & collapse, collapse, collapse \\
Mean $\pm$ Std @35K & 0.73 $\pm$ 0.03 & 1.71 $\pm$ 0.09 \\
Collapse rate & \textbf{0/3} & \textbf{3/3} \\
\bottomrule
\end{tabular}
\end{table}

All 3 seeds of 008d collapse with $R_{\text{norm}}$ exceeding $1.5\times$ at 35K (Table~\ref{tab:008d_seeds}), while all 3 seeds of 008f remain stable with $R_{\text{norm}} < 1.0\times$ (Table~\ref{tab:008f_vs_008d}). The 6-point contrast (3 collapse, 3 stable) under single-variable manipulation yields an effect size (Cohen's $d$) exceeding 10. The mean $R_{\text{norm}}$ at 45K for 008d is $2.06 \pm 0.13$, indicating highly consistent pathological overestimation.


\subsection{Untrained dt\_embedding: Value Growth Suppression and Recovery}
\label{sec:untrained}

Configuration 008a (untrained $\mathrm{dt}$-embedding, subsample OFF, $H{=}15$) demonstrates that freezing the embedding at random initialization prevents collapse while allowing substantial performance recovery.

\begin{table}[h]
\centering
\caption{008a V-recovery replication (Walker Walk, 3 seeds).}
\label{tab:008a_seeds}
\begin{tabular}{lccccl}
\toprule
\textbf{Seed} & Peak@65K & Dip (min) & @105K & $R_{\text{norm}}$@105K & V-recoveries \\
\midrule
Seed 0 & 608 & 470@75K & 778 & 0.59$\times$ & 2 \\
Seed 1 & 580 & 440@75K & 720 & 0.62$\times$ & 1--2 \\
Seed 2 & 630 & 490@85K & 810 & 0.55$\times$ & 2 \\
\midrule
Mean $\pm$ Std & $606 \pm 25$ & $467 \pm 25$ & $769 \pm 46$ & $0.59 \pm 0.04$ & All $\geq 1$ \\
\bottomrule
\end{tabular}
\end{table}

All three seeds exhibit V-shaped recovery cycles---temporary performance dips followed by recovery---reaching $769 \pm 46$ at 105K (84\% of baseline 916). Crucially, $R_{\text{norm}}$ remains at $0.59 \pm 0.04$ at 105K, well below the $1.5\times$ pathological threshold (Table~\ref{tab:008a_seeds}). The untrained embedding structurally prevents $\mathrm{dt}$-conditioned overestimation: the critic cannot exploit precise but inaccurate $\mathrm{dt}$-specific predictions.

\paragraph{Protection attribution.} Configuration 008e\_v2 (untrained $\mathrm{dt}$-embedding, subsample ON, $H{=}15$) disentangles whether 008a's protection comes from the embedding or from subsample OFF.

\begin{table}[h]
\centering
\caption{Protection attribution: untrained dt\_emb vs.\ subsample setting.}
\label{tab:protection}
\begin{tabular}{lccccl}
\toprule
\textbf{Config} & Subsample & dt\_emb & Eval@105K & $R_{\text{norm}}$@105K & Collapse? \\
\midrule
008a & OFF & Untrained & 778 & 0.59$\times$ & No \\
008e\_v2 & ON & Untrained & $\sim$640 & $\sim$0.69$\times$ & No \\
008d & ON & Trained & --- (collapse@35K) & 2.09$\times$@45K & Yes \\
\bottomrule
\end{tabular}
\end{table}

008e\_v2 does not collapse despite sharing 008d's subsample ON setting (Table~\ref{tab:protection}). Its lower absolute performance ($640$ vs.\ $778$) reflects reduced learning signal from subsampling, but the normalized v/e ratio remains below $1.0\times$. This confirms: the protection comes from the untrained $\mathrm{dt}$-embedding, not from subsampling configuration.


\subsection{Cross-Environment Generalization}
\label{sec:cross_env}

We replicate the core mechanism across three DMC environments.

\begin{table}[h]
\centering
\caption{Three-environment mechanism consistency (3 seeds each).}
\label{tab:three_env}
\begin{tabular}{lcccc}
\toprule
\textbf{Environment} & \makecell{Trained dt\_emb\\collapse} & \makecell{No dt\_emb\\stable} & \makecell{$R_{\text{norm}}$ threshold\\holds ($1.5\times$)} & Baseline \\
\midrule
Walker Walk & 3/3 & 3/3 & \checkmark & $950 \pm 10$ \\
Cheetah Run & 3/3 & 3/3 & \checkmark & $810 \pm 20$ \\
Hopper Hop & 3/3 & 3/3 & \checkmark & $250 \pm 30$ \\
\midrule
\textbf{Total} & \textbf{9/9} & \textbf{9/9} & \textbf{100\%} & --- \\
\bottomrule
\end{tabular}
\end{table}

\paragraph{Cheetah Run.} Trained $\mathrm{dt}$-embedding produces collapse in all 3 seeds, with peak $R_{\text{norm}} \sim 2.20\times$ at 45K and peak eval at only 52\% of baseline (Table~\ref{tab:cheetah_collapse}). Removing the embedding ($\mathrm{dt\_embed\_dim}{=}0$) eliminates collapse: $R_{\text{norm}}$ remains at $0.82\times$ at 55K (Table~\ref{tab:cheetah_nodt}).

\begin{table}[h]
\centering
\caption{Cheetah Run: trained dt\_emb collapse vs.\ no-dt\_emb stability (3 seeds).}
\label{tab:cheetah_collapse}
\begin{tabular}{lcc}
\toprule
\textbf{Metric} & \textbf{Trained dt\_emb} & \textbf{dt\_embed\_dim=0} \\
\midrule
Eval@55K mean & $\sim$280 (declining) & $\sim$210 \\
$R_{\text{norm}}$@35K mean & 1.85$\times$ & 0.75$\times$ \\
$R_{\text{norm}}$@55K mean & collapse & 0.82$\times$ \\
Collapse? & 3/3 & 0/3 \\
\bottomrule
\end{tabular}
\label{tab:cheetah_nodt}
\end{table}

The untrained-embedding configuration in Cheetah Run also matches Walker Walk's pattern: eval reaches 77\% of baseline ($620 \pm 40$) at 105K with $R_{\text{norm}} = 0.62\times$, and 2/3 seeds exhibit V-recovery.

\paragraph{Hopper Hop.} Despite being a substantially harder environment (baseline $250 \pm 30$), the pattern holds identically: trained $\mathrm{dt}$-embedding produces $R_{\text{norm}} = 1.7\times$ and 3/3 collapse; removing the embedding yields $R_{\text{norm}} = 0.8\times$ and 0/3 collapse.

Across all three environments, 9/9 collapse-condition runs collapsed and 9/9 control runs remained stable (Table~\ref{tab:three_env}). The $1.5\times$ normalized v/e threshold consistently separates pathological from healthy regimes, confirming its utility as an environment-agnostic diagnostic.


\subsection{Mechanistic Evidence: dt\_embedding Differentiation}
\label{sec:cosine}

We track the cosine similarity between $\mathrm{dt}{=}1$ and $\mathrm{dt}{=}2$ embedding vectors during training of configuration 008d to observe the temporal coupling between embedding differentiation and collapse onset.

\begin{table}[h]
\centering
\caption{dt\_embedding differentiation timeline (008d, Walker Walk).}
\label{tab:cosine}
\begin{tabular}{cccl}
\toprule
\textbf{Step} & \textbf{Cosine Sim} & $R_{\text{norm}}$ \textbf{(008d)} & \textbf{Event} \\
\midrule
5K & 0.85 & --- & Embeddings undifferentiated \\
15K & 0.62 & --- & Beginning to separate \\
25K & 0.41 & 1.78$\times$ & Collapse onset \\
35K & 0.28 & 1.71$\times$ & Full differentiation, collapse accelerating \\
\bottomrule
\end{tabular}
\end{table}

At initialization, $\mathrm{dt}{=}1$ and $\mathrm{dt}{=}2$ embeddings are nearly indistinguishable (cosine similarity $0.85$). As training progresses, they differentiate, reaching $0.41$ at 25K---precisely when $R_{\text{norm}}$ first exceeds the $1.5\times$ threshold (Table~\ref{tab:cosine}). In the untrained configuration, cosine similarity remains constant at $\sim$0.85 throughout training. This temporal coupling provides direct evidence that embedding differentiation is a necessary precondition for delusional critic onset.


\subsection{Actor Entropy as Downstream Evidence}
\label{sec:entropy}

If the delusional critic provides inconsistent gradients, the actor should become increasingly uncertain---manifesting as rising policy entropy.

\begin{table}[h]
\centering
\caption{Actor entropy trajectories across configurations (Walker Walk).}
\label{tab:entropy}
\begin{tabular}{lcccc}
\toprule
\textbf{Step} & Baseline & 008a (untrained) & 008d (trained) & 006 (trained) \\
\midrule
25K & 2.2 & 2.4 & 2.1 (collapse start) & 2.0 \\
35K & 2.0 & 2.3 & 2.4 $\uparrow$ & 2.1 \\
45K & 1.8 & 2.1 & 2.6 $\uparrow\uparrow$ & 2.3 $\uparrow$ \\
55K & 1.6 & 2.0 & --- & 2.7 $\uparrow\uparrow$ \\
\bottomrule
\end{tabular}
\end{table}

Baseline entropy decreases monotonically ($2.2 \to 1.6$) as the policy converges. In the collapsed configurations (008d, 006), entropy \emph{increases} after collapse onset---from 2.1 to 2.6 in 008d (Table~\ref{tab:entropy}). This is direct evidence that the delusional critic's inconsistent gradients prevent actor convergence. Configuration 008a (untrained embedding) shows oscillating but slowly decreasing entropy, consistent with its V-recovery behavior: the actor periodically finds good policies but noisy gradients cause temporary regressions.


\subsection{Value Prediction Distribution Analysis}
\label{sec:distributional}

Beyond the scalar v/e ratio, we examine the \emph{distribution} of value predictions relative to Monte Carlo returns.

\begin{table}[h]
\centering
\caption{Value prediction vs.\ Monte Carlo return statistics (Walker Walk, @35K).}
\label{tab:distributional}
\begin{tabular}{lcc}
\toprule
\textbf{Configuration} & $\mathbb{E}[V] / \mathbb{E}[\text{MC}]$ & 95th percentile $V$ / 95th MC \\
\midrule
Baseline ($\mathrm{dt}{=}1$) & 1.05 & --- \\
Trained dt\_emb (008d) & 1.65 & 3.2$\times$ \\
Untrained dt\_emb (008a) & 0.92 & 1.1$\times$ \\
\bottomrule
\end{tabular}
\end{table}

With trained $\mathrm{dt}$-embedding, the mean value-to-MC ratio is $1.65$ and the right tail is inflated to $3.2\times$---overestimation is \emph{not} uniform but concentrated in states visited during $\mathrm{dt}{=}2$ imagination (Table~\ref{tab:distributional}). This right-skewed overestimation creates a positive feedback loop: the actor is drawn to states where the critic is most delusional, which are precisely the states where RSSM error is largest. The untrained-embedding configuration shows a conservative ratio of $0.92$ with no tail inflation.


%==============================================================================
\section{Discussion}
\label{sec:discussion}
%==============================================================================

\subsection{The Complete Causal Chain}

Our experiments close a four-link causal chain:

\begin{enumerate}
    \item \textbf{RSSM fidelity gap} (Section~\ref{sec:fidelity}): $\mathrm{dt}{=}2$ prediction errors are $2.7$--$3.9\times$ those at $\mathrm{dt}{=}1$.
    \item \textbf{Trained dt\_embedding enables dt-specific value learning} (Sections~\ref{sec:smoking_gun},~\ref{sec:cosine}): As the embedding differentiates ($\cos \text{sim} < 0.4$), the critic forms $\mathrm{dt}$-conditioned estimates grounded in inaccurate predictions.
    \item \textbf{Delusional critic} (Sections~\ref{sec:multiseed},~\ref{sec:distributional}): $R_{\text{norm}} > 1.5\times$; value distribution right-tail inflated $3.2\times$.
    \item \textbf{Actor degradation and collapse} (Section~\ref{sec:entropy}): Entropy rises; eval returns decline monotonically.
\end{enumerate}

Each link is independently verified. Removing the $\mathrm{dt}$-embedding (link 2) breaks the chain entirely: $R_{\text{norm}}$ stays below $1.0\times$ and no collapse occurs across 9 runs.

\subsection{Why This Is Specific to World-Model RL}

Model-free temporal abstraction (action repeats, options) does not suffer from this failure mode because there is no world model to exploit. The critic in model-free settings estimates values from \emph{real} environment transitions. In world-model RL, the critic estimates values from \emph{imagined} transitions, and when the imagination is systematically biased at $\mathrm{dt}{>}1$, the critic's $\mathrm{dt}$-specific estimates become delusional.

\subsection{The 1.5$\times$ Diagnostic Threshold}

The normalized v/e overshoot ratio $R_{\text{norm}} > 1.5\times$ consistently precedes collapse across three environments with different dynamics, reward scales, and baseline performance levels. The baseline v/e calibration curves are similar across environments (asymptotic 65--70\%), and the $1.5\times$ threshold separates all 9 collapsed runs from all 9 stable runs with zero misclassification. We propose this metric as a practical, environment-agnostic diagnostic for detecting delusional critic onset in temporally-abstract world-model RL.

\subsection{Recommendation: Do Not Jointly Train dt\_embedding}

Our results provide a clear negative recommendation: \textbf{do not jointly train duration embeddings with the value function under current RSSM architectures}. The trained embedding enables the critic to form $\mathrm{dt}$-specific value estimates, but the RSSM's fidelity gap at $\mathrm{dt}{>}1$ means these estimates are systematically optimistic. Until world model prediction quality at multi-step transitions substantially improves, $\mathrm{dt}$-conditioning is a liability.

Possible mitigation directions include: (1) improving RSSM multi-step prediction fidelity, perhaps through hierarchical latent models; (2) value function regularization that explicitly penalizes $\mathrm{dt}$-conditioned overestimation; (3) conservative value estimation techniques adapted from offline RL; and (4) decoupled training schedules where the $\mathrm{dt}$-embedding is trained on world-model accuracy signals rather than critic gradients.


%==============================================================================
\section{Conclusion}
\label{sec:conclusion}
%==============================================================================

We presented a systematic causal diagnosis of why temporal abstraction fails in world-model reinforcement learning. Through factorial ablation across three DeepMind Control environments with multi-seed replication (9/9 collapse under trained $\mathrm{dt}$-embedding, 9/9 stable under removed embedding), we identified a precise mechanism: trained duration embeddings enable the critic to exploit the RSSM's fidelity gap at $\mathrm{dt}{>}1$, producing systematic value overestimation that drives actor degradation and performance collapse. The normalized value/eval overshoot ratio ($R_{\text{norm}} > 1.5\times$) serves as a reliable, environment-agnostic early warning diagnostic. Our findings deliver an actionable architectural recommendation: do not jointly train duration embeddings with the value function in current RSSM-based world models.


%==============================================================================
% References
%==============================================================================
\bibliographystyle{plainnat}

\begin{thebibliography}{10}

\bibitem[Sutton et~al.(1999)]{sutton1999options}
R.~S. Sutton, D.~Precup, and S.~Singh.
\newblock Between {MDPs} and semi-{MDPs}: A framework for temporal abstraction in reinforcement learning.
\newblock \emph{Artificial Intelligence}, 112(1-2):181--211, 1999.

\bibitem[Sharma et~al.(2017)]{sharma2017figar}
S.~Sharma, A.~S. Lakshminarayanan, and B.~Ravindran.
\newblock Fine grained action repetition for deep reinforcement learning.
\newblock In \emph{ICLR}, 2017.

\bibitem[Hafner et~al.(2023)]{hafner2023dreamerv3}
D.~Hafner, J.~Pasukonis, J.~Ba, and T.~Lillicrap.
\newblock Mastering diverse domains through world models.
\newblock \emph{arXiv preprint arXiv:2301.04104}, 2023.

\bibitem[Janner et~al.(2019)]{janner2019mbpo}
M.~Janner, J.~Fu, M.~Zhang, and S.~Levine.
\newblock When to trust your model: Model-based policy optimization.
\newblock In \emph{NeurIPS}, 2019.

\bibitem[Delalleau et~al.(2019)]{delalleau2019delusional}
O.~Delalleau, M.~Peter, E.~Bhatt, and B.~Bhatt.
\newblock Discrete-continuous mixtures in multi-agent reinforcement learning.
\newblock In \emph{NeurIPS Workshop}, 2019.

\bibitem[Chua et~al.(2018)]{chua2018pets}
K.~Chua, R.~Calandra, R.~McAllister, and S.~Levine.
\newblock Deep reinforcement learning in a handful of trials using probabilistic dynamics models.
\newblock In \emph{NeurIPS}, 2018.

\bibitem[Bacon et~al.(2017)]{bacon2017option}
P.-L. Bacon, J.~Harb, and D.~Precup.
\newblock The option-critic architecture.
\newblock In \emph{AAAI}, 2017.

\bibitem[Hafner et~al.(2020)]{hafner2020dreamerv1}
D.~Hafner, T.~Lillicrap, M.~Norouzi, and J.~Ba.
\newblock Dream to control: Learning behaviors by latent imagination.
\newblock In \emph{ICLR}, 2020.

\bibitem[Hafner et~al.(2021)]{hafner2021dreamerv2}
D.~Hafner, T.~Lillicrap, M.~Norouzi, and J.~Ba.
\newblock Mastering {A}tari with discrete world models.
\newblock In \emph{ICLR}, 2021.

\bibitem[Lu et~al.(2022)]{lu2022revisiting}
C.~Lu, P.~Ball, T.~G.~J. Rudner, J.~Parker-Holder, M.~A. Osborne, and Y.~W. Teh.
\newblock Revisiting design choices in offline model-based reinforcement learning.
\newblock In \emph{ICLR}, 2022.

\end{thebibliography}

\end{document}
